\documentclass[10pt,a4paper]{amsart}
\usepackage{amsmath,amssymb,amsthm,graphicx}
% Retained minimal preamble for the Stacks Project chapter sites-cohomology.tex (Cohomology on Sites).
% Verbatim excerpt of the chapter preamble: Xy-pic setup (native preamble lines 14-18) and the native
% theorem declarations with the mathematical macros (native preamble lines 154-217). Copied mechanically;
% nothing added, removed or redefined. The wrapper class and the amsmath/amssymb/amsthm/graphicx base
% packages are supplied by the assembler, so the class/externaldocument part of the native preamble is omitted.
% For commutative diagrams we use Xy-pic
\usepackage[all]{xy}

% We use 2cell for 2-commutative diagrams.
\xyoption{2cell}
% Theorem environments.
%
\theoremstyle{plain}
\newtheorem{theorem}[subsection]{Theorem}
\newtheorem{proposition}[subsection]{Proposition}
\newtheorem{lemma}[subsection]{Lemma}

\theoremstyle{definition}
\newtheorem{definition}[subsection]{Definition}
\newtheorem{example}[subsection]{Example}
\newtheorem{exercise}[subsection]{Exercise}
\newtheorem{situation}[subsection]{Situation}

\theoremstyle{remark}
\newtheorem{remark}[subsection]{Remark}
\newtheorem{remarks}[subsection]{Remarks}

\numberwithin{equation}{subsection}

% Macros
%
\def\lim{\mathop{\mathrm{lim}}\nolimits}
\def\colim{\mathop{\mathrm{colim}}\nolimits}
\def\Spec{\mathop{\mathrm{Spec}}}
\def\Hom{\mathop{\mathrm{Hom}}\nolimits}
\def\Ext{\mathop{\mathrm{Ext}}\nolimits}
\def\SheafHom{\mathop{\mathcal{H}\!\mathit{om}}\nolimits}
\def\SheafExt{\mathop{\mathcal{E}\!\mathit{xt}}\nolimits}
\def\Sch{\mathit{Sch}}
\def\Mor{\mathop{\mathrm{Mor}}\nolimits}
\def\Ob{\mathop{\mathrm{Ob}}\nolimits}
\def\Sh{\mathop{\mathit{Sh}}\nolimits}
\def\NL{\mathop{N\!L}\nolimits}
\def\CH{\mathop{\mathrm{CH}}\nolimits}
\def\proetale{{pro\text{-}\acute{e}tale}}
\def\etale{{\acute{e}tale}}
\def\QCoh{\mathit{QCoh}}
\def\Ker{\mathop{\mathrm{Ker}}}
\def\Im{\mathop{\mathrm{Im}}}
\def\Coker{\mathop{\mathrm{Coker}}}
\def\Coim{\mathop{\mathrm{Coim}}}

% Boxtimes
%
\DeclareMathSymbol{\boxtimes}{\mathbin}{AMSa}{"02}

%
% Macros for moduli stacks/spaces
%
\def\QCohstack{\mathcal{QC}\!\mathit{oh}}
\def\Cohstack{\mathcal{C}\!\mathit{oh}}
\def\Spacesstack{\mathcal{S}\!\mathit{paces}}
\def\Quotfunctor{\mathrm{Quot}}
\def\Hilbfunctor{\mathrm{Hilb}}
\def\Curvesstack{\mathcal{C}\!\mathit{urves}}
\def\Polarizedstack{\mathcal{P}\!\mathit{olarized}}
\def\Complexesstack{\mathcal{C}\!\mathit{omplexes}}
% \Pic is the operator that assigns to X its picard group, usage \Pic(X)
% \Picardstack_{X/B} denotes the Picard stack of X over B
% \Picardfunctor_{X/B} denotes the Picard functor of X over B
\def\Pic{\mathop{\mathrm{Pic}}\nolimits}
\def\Picardstack{\mathcal{P}\!\mathit{ic}}
\def\Picardfunctor{\mathrm{Pic}}
\def\Deformationcategory{\mathcal{D}\!\mathit{ef}}

% --- appended 2026-10-12 (Chief): the project own definitions of the reference/slogan/history
% display environments, copied mechanically from the shared Stacks preamble lines 8-12
% (\usepackage{verbatim} plus the three \newenvironment lines, which the chapter preamble keeps
% verbatim from the project preamble). Nothing added, removed or redefined: these environments
% carry no mathematical content and the project itself typesets them as comments. Without these
% three lines a counted statement whose own text uses them fails with Environment slogan undefined.
\usepackage{verbatim}
\newenvironment{reference}{\comment}{\endcomment}
\newenvironment{slogan}{\comment}{\endcomment}
\newenvironment{history}{\comment}{\endcomment}
\usepackage{cleveref}
\crefname{originallemma}{original lemma}{original lemmas}
\Crefname{originallemma}{Original lemma}{Original lemmas}
\makeatletter
\newlabel{lemma-olsson-laszlo}{{0D6S}{source}{}{sourcepointer-1}{}}
\newlabel{lemma-olsson-laszlo@cref}{{[originallemma][0][]0D6S}{[1][0][]source}{}{}{}}
\makeatother
\crefname{originallemma}{original lemma}{original lemmas}
\Crefname{originallemma}{Original lemma}{Original lemmas}
\makeatletter
\newlabel{lemma-Rf-commutes-with-Rlim}{{0A07}{source}{}{sourcepointer-2}{}}
\newlabel{lemma-Rf-commutes-with-Rlim@cref}{{[originallemma][0][]0A07}{[1][0][]source}{}{}{}}
\makeatother
\crefname{originallemma}{original lemma}{original lemmas}
\Crefname{originallemma}{Original lemma}{Original lemmas}
\makeatletter
\newlabel{derived-lemma-two-ss-complex-functor}{{015J}{source}{}{sourcepointer-3}{}}
\newlabel{derived-lemma-two-ss-complex-functor@cref}{{[originallemma][0][]015J}{[1][0][]source}{}{}{}}
\makeatother
\crefname{originallemma}{original lemma}{original lemmas}
\Crefname{originallemma}{Original lemma}{Original lemmas}
\makeatletter
\newlabel{homology-lemma-biregular-ss-converges}{{012W}{source}{}{sourcepointer-4}{}}
\newlabel{homology-lemma-biregular-ss-converges@cref}{{[originallemma][0][]012W}{[1][0][]source}{}{}{}}
\makeatother
\crefname{originallemma}{original lemma}{original lemmas}
\Crefname{originallemma}{Original lemma}{Original lemmas}
\makeatletter
\newlabel{lemma-olsson-laszlo-modified}{{0D6T}{source}{}{sourcepointer-5}{}}
\newlabel{lemma-olsson-laszlo-modified@cref}{{[originallemma][0][]0D6T}{[1][0][]source}{}{}{}}
\makeatother
\title{A version of the Olsson-Laszlo boundedness lemma for a morphism of ringed topoi}
\author{The Stacks Project Authors}
\date{stacks-project:sites-cohomology.tex:local-snapshot-20261011; GNU FDL 1.2 or later}
\begin{document}
\maketitle
\noindent\textit{Editorial note.} All delivered mathematical text is the author's own, copied from the pinned Stacks Project chapter file sites-cohomology.tex; no statement or proof step was rewritten, reordered or removed.\par\medskip
\noindent\textit{Editorial note.} The proof is the author's complete argument for both assertions of the statement: the derived limit presentation, the placement of each truncated pushforward in the bounded below derived category and the degree comparison are all written out, with every cited input named.\par\medskip
\noindent\textit{Editorial note.} Every labelled reference occurring inside the counted spans is carried in assembly.external\_references as a non-load-bearing navigation pointer (same-chapter ones at their definition site, cross-chapter ones at the invocation inside the counted span), so no load-bearing mathematics is replaced by a stub. No bibliography citation occurs inside the counted spans: the author's statement environment carries, inside a 'reference' environment at source lines 5797-5800, a provenance remark citing a published lemma of Olsson and Laszlo, and that remark is not part of the mathematical statement; it is therefore not delivered, and the counted statement is split into the two spans 5795-5796 and 5801-5823. No citation was supplied, invented or reconstructed, and no other text of the author was dropped.\par\medskip
\noindent\textit{Editorial note.} The checked chapter list for same-name collisions is the set of stacks chapters staged in this repository: chow, cohomology, divisors, etale-cohomology, flat, more-algebra, more-morphisms, morphisms, spaces-more-morphisms, varieties together with this card's chapter sites-cohomology.\par\medskip
\noindent\textit{Editorial note.} The preamble is the chapter's own retained minimal preamble (sources/stacks-sites-cohomology/retained-preamble.tex), a verbatim excerpt of the native preamble: the Xy-pic setup and the native theorem declarations with the mathematical macros.\par\medskip
\noindent\textit{Allowed standard prerequisites (editorial).}\begin{itemize}
\item The Olsson-Laszlo situation. Same chapter, Situation situation-olsson-laszlo. Supplies the hypotheses (3) and (4) refer to for both pairs (C, O, A) and (C', O', A').
\item Unbounded complexes are derived limits of their truncations. Same chapter, Lemma lemma-olsson-laszlo. Writes K as the derived limit of its canonical truncations.
\item The derived pushforward commutes with derived limits. Same chapter, Lemma lemma-Rf-commutes-with-Rlim. Moves the derived limit through Rf\_*.
\item The spectral sequence of a complex with respect to a functor. Derived Categories, Lemma derived-lemma-two-ss-complex-functor, together with Homology, Lemma homology-lemma-biregular-ss-converges. Shows that each truncated pushforward lies in the bounded below derived category of A'.
\item The modified Olsson-Laszlo comparison. Same chapter, Lemma lemma-olsson-laszlo-modified. Converts the comparison of truncations into the two displayed assertions.
\end{itemize}\medskip
\ifdefined\hypertarget\hypertarget{sourcepointer-1}{}\fi
\noindent\textit{Source reference.} Original lemma 0D6S: sources/stacks-sites-cohomology/pinned-original-source.tex, source lines 5707--5707. Navigation pointer invoked at proof line 5826; the counted span contains the pointer only and the cited result is not reproduced inside it, so the pointer is registered as a non-load-bearing navigation stub in the established form for Stacks material. No bibliography citation is involved.\par\medskip
\ifdefined\hypertarget\hypertarget{sourcepointer-2}{}\fi
\noindent\textit{Source reference.} Original lemma 0A07: sources/stacks-sites-cohomology/pinned-original-source.tex, source lines 5150--5150. Navigation pointer invoked at proof line 5827; the counted span contains the pointer only and the cited result is not reproduced inside it, so the pointer is registered as a non-load-bearing navigation stub in the established form for Stacks material. No bibliography citation is involved.\par\medskip
\ifdefined\hypertarget\hypertarget{sourcepointer-3}{}\fi
\noindent\textit{Source reference.} Original lemma 015J: sources/stacks-sites-cohomology/pinned-original-source.tex, source lines 5742--5742. Navigation pointer invoked at proof line 5742; the counted span contains the pointer only and the cited result is not reproduced inside it, so the pointer is registered as a non-load-bearing navigation stub in the established form for Stacks material. No bibliography citation is involved.\par\medskip
\ifdefined\hypertarget\hypertarget{sourcepointer-4}{}\fi
\noindent\textit{Source reference.} Original lemma 012W: sources/stacks-sites-cohomology/pinned-original-source.tex, source lines 5838--5838. Navigation pointer invoked at proof line 5838; the counted span contains the pointer only and the cited result is not reproduced inside it, so the pointer is registered as a non-load-bearing navigation stub in the established form for Stacks material. No bibliography citation is involved.\par\medskip
\ifdefined\hypertarget\hypertarget{sourcepointer-5}{}\fi
\noindent\textit{Source reference.} Original lemma 0D6T: sources/stacks-sites-cohomology/pinned-original-source.tex, source lines 5726--5726. Navigation pointer invoked at proof line 5845; the counted span contains the pointer only and the cited result is not reproduced inside it, so the pointer is registered as a non-load-bearing navigation stub in the established form for Stacks material. No bibliography citation is involved.\par\medskip
\begin{situation}
\label{situation-olsson-laszlo}
Let $\mathcal{C}$ be a site. Let $\mathcal{O}$ be a sheaf of rings on
$\mathcal{C}$. Let $\mathcal{A} \subset \textit{Mod}(\mathcal{O})$
be a weak Serre subcategory. We assume the following is true:
there exists a subset $\mathcal{B} \subset \Ob(\mathcal{C})$ such that
\begin{enumerate}
\item every object of $\mathcal{C}$ has a covering whose
members are in $\mathcal{B}$, and
\item for every $V \in \mathcal{B}$ there exists an integer $d_V$
and a cofinal system $\text{Cov}_V$ of coverings of $V$ such
that
$$
H^p(V_i, \mathcal{F}) = 0 \text{ for }
\{V_i \to V\} \in \text{Cov}_V,\ p > d_V, \text{ and }
\mathcal{F} \in \Ob(\mathcal{A})
$$
\end{enumerate}
\end{situation}

\begin{lemma}
\label{lemma-olsson-laszlo-map-version-one}

Let $f : (\Sh(\mathcal{C}), \mathcal{O}) \to (\Sh(\mathcal{C}'), \mathcal{O}')$
be a morphism of ringed topoi.
Let $\mathcal{A} \subset \textit{Mod}(\mathcal{O})$
and $\mathcal{A}' \subset \textit{Mod}(\mathcal{O}')$
be weak Serre subcategories. Assume there is an integer $N$ such that
\begin{enumerate}
\item $\mathcal{C}, \mathcal{O}, \mathcal{A}$ satisfy the
assumption of Situation \ref{situation-olsson-laszlo},
\item $\mathcal{C}', \mathcal{O}', \mathcal{A}'$ satisfy the
assumption of Situation \ref{situation-olsson-laszlo},
\item $R^pf_*\mathcal{F} \in \Ob(\mathcal{A}')$ for
$p \geq 0$ and $\mathcal{F} \in \Ob(\mathcal{A})$,
\item $R^pf_*\mathcal{F} = 0$ for
$p > N$ and $\mathcal{F} \in \Ob(\mathcal{A})$,
\end{enumerate}
Then for $K$ in $D_\mathcal{A}(\mathcal{O})$ we have
\begin{enumerate}
\item[(a)] $Rf_*K$ is in $D_{\mathcal{A}'}(\mathcal{O}')$,
\item[(b)] the map
$H^j(Rf_*K) \to H^j(Rf_*(\tau_{\geq -n}K))$ is an isomorphism
for $j \geq N - n$.
\end{enumerate}
\end{lemma}

\begin{proof}
By Lemma \ref{lemma-olsson-laszlo} we have $K = R\lim \tau_{\geq -n}K$.
By Lemma \ref{lemma-Rf-commutes-with-Rlim}
we have $Rf_*K = R\lim Rf_*\tau_{\geq -n}K$.
The complexes $Rf_*\tau_{\geq -n}K$ are bounded below.
The spectral sequence
$$
E_2^{p, q} = R^pf_*H^q(\tau_{\geq -n}K)
$$
converging to $H^{p + q}(Rf_*\tau_{\geq -n}K)$
(Derived Categories, Lemma \ref{derived-lemma-two-ss-complex-functor})
and assumption (3)
show that $Rf_*\tau_{\geq -n}K$ lies in $D^+_{\mathcal{A}'}(\mathcal{O}')$,
see Homology, Lemma \ref{homology-lemma-biregular-ss-converges}.
Observe that for $m \geq n$ the map
$$
Rf_*(\tau_{\geq -m}K) \longrightarrow Rf_*(\tau_{\geq -n}K)
$$
induces an isomorphism on cohomology sheaves in degrees $j \geq -n + N$
by the spectral sequences above. Hence we may apply
Lemma \ref{lemma-olsson-laszlo-modified} to conclude.
\end{proof}

\end{document}
